\hypertarget{sec:corporate-residuals}{%
\subsection{Changed obligations and unfinished production}\label{sec:corporate-residuals}}

A company can change its restoration estimate before a machine wears out, promise to repair sold products, or end a period with unfinished goods. Each condition changes a different part of its accounts. The following cases extend the preceding operating methods through those transitions. They retain exact illustrative amounts and separate economic evidence from bookkeeping.

Restoration, production and uncertain-tax methods use the retained historical EU text; assurance and mixed contracts use the inspected AASB compilations for synthetic 2024/2025 Australian for-profit reporters. The AASB 137 PDF covers annual periods beginning from 1 January 2023 to before 1 July 2026; its website gives a different end-period summary. Both admit these particular historical calendar years. A real application still needs its own source and period admission. Amounts are pretax unless an explicit tax calculation follows.

\hypertarget{restoration-estimates-change-the-asset-and-the-obligation-differently}{%
\subsubsection{Restoration estimates change the asset and the obligation differently}\label{restoration-estimates-change-the-asset-and-the-obligation-differently}}

The restoration record retains the related component and liability, cost/revaluation policy, gross cost/accumulated depreciation, remaining life/residual, current cash-flow estimate and rate, prior estimate identity, same-asset revaluation surplus/prior profit deficit, hypothetical cost-model carrying amount, impairment evidence and whole-class valuation review. IFRIC~1.2 \citep[IFRIC~1.1--8]{eu20231803} requires the original liability to have been included in the asset's cost. A standalone litigation obligation is not admitted to this asset-adjustment method.

Separate passage of time from a new estimate. Under the cost model, a liability increase adds to PPE; a decrease reduces PPE only to zero, with excess in profit. An increase prompts assessment of impairment indications and an IAS36 test when required. Depreciate the adjusted carrying amount less admitted residual over the remaining useful life. Discount unwind is finance expense, never capitalised under IFRIC1.8. Once useful life is exhausted, later estimate changes go to profit under7 regardless of asset model. Actual settlement reduces the obligation and cash; a difference from the last carrying amount requires a dated remeasurement.

With cash100, PPE80 and provision20, opening equity is160. An admitted new provision50 adds30 to PPE. Recoverable amount90 requires impairment20: PPE90, cash100, liability50 and equity140. Four years remaining and zero residual gives annual depreciation22.5. The subsequent depreciation-only illustration assumes a separately supplied zero discount rate and no further liability change; PPE67.5 and equity117.5 follow. In another branch, PPE30/provision80 and a new provision30 require Dr provision50 / Cr PPE30 / Cr profit20. Negative PPE20 would be wrong. Cash100 then equals liability30 plus equity70.

Under the revaluation model, a liability increase first consumes a same-asset surplus through OCI, then affects profit; a decrease reverses a same-asset profit deficit and otherwise increases OCI. A decrease exceeding the hypothetical cost-model carrying amount has an immediate profit component. Determine any necessary class revaluation first so that these splits use the correct history. With PPE150/provision50/surplus20, an increase35 gives OCI loss20 and profit loss15; PPE stays150 only because a supplied class valuation review supports it. Cash100 then gives equity165 and liability85. Surplus on another asset cannot be borrowed. In a separate decrease case, PPE25, hypothetical cost40, earlier profit deficit15 and provision80→30 yield excess10 plus deficit reversal15 in profit and25 in OCI. Cash100+PPE25=provision30+equity95. Using25 as the cost floor would compute the wrong split. The prior deficit explains why25 and40 differ; no fresh impairment is inferred from that difference.

A full cash history opens cash100/PPE120/provision20/equity200. Two-year straight-line life and zero residual give depreciation60 in year1. At10\%, unwind2 raises the provision22; new supported estimate33 adds11 to PPE. Close PPE71/provision33/equity138, with pretax loss62 and no cash change. Year2 depreciation71 and unwind3.3 give loss74.3; actual settlement36.3 leaves cash/equity63.7 and zero PPE/provision. Across both years depreciation131 plus finance5.3 equals cumulative loss136.3. The total cash bridge is exact; IAS7 classification of restoration settlement needs a separate fact-specific policy. The method therefore returns a cash-classification obligation rather than inventing a complete categorized cash-flow statement. A fully exhausted PPE0 with provision20→25 gives expense5 even if a historic revaluation reserve exists.

\hypertarget{assurance-is-a-measured-obligation-alongside-the-product-sale}{%
\subsubsection{Assurance is a measured obligation alongside the product sale}\label{assurance-is-a-measured-obligation-alongside-the-product-sale}}

The assurance record stores covered sales/population, legal or constructive promise, coverage/task findings, claim history, remaining-outcome law, timing/discount, reimbursement contract and virtual-certainty evidence. AASB15.B28--B33 determines whether assurance, service or an inseparable combined service is present. AASB~137.14/24 \citep[paras.~14,24,36--62]{aasb137dated} admission precedes measurement. For a large population39 uses weighted outcomes; individual claim probability does not replace the population outflow judgment. Correlation matters for distributions even when marginal means are sufficient to add expected cost. A single lawsuit uses40's different analysis.

Consider100 products with admitted per-product costs0,10,40 and probabilities 3/4,1/5,1/20. Expected cost is100×(0+2+2)=400. The entire population has a present assurance duty, probable outflow and reliable estimate; discount is immaterial over the short horizon. Opening cash1000 and inventory600 are funded by equity1600. Cash sales2000 and inventory expense600 accompany warranty expense/provision400. An independently enforceable quarter-cost reimbursement is virtually certain on settlement, so recognise a separate asset100/benefit100. Do not net it against the provision. Paragraph54 permits net expense presentation.

Pay claims100 and receive reimbursement25. A new remaining-claims estimate250 replaces the old300, releasing50; reduce the remaining reimbursement75 to62.5, expense12.5. Year1 cash2925+reimbursement62.5=provision250+pretax equity2737.5. Profit1137.5 follows from2000−600−400+100+50−12.5. Direct operating cash1925 also equals profit1137.5+inventory release600+provision increase250−receivable increase62.5. Year2 settlement250 and receipt62.5 leave cash/equity2737.5 with no further profit under the realised amounts. Without virtual certainty, no reimbursement asset is recognised; retain the separate uncertainty. Missing claim evidence does not permit zero provision or the printed400 by default.

\hypertarget{a-mixed-amendment-needs-two-revenue-calculations}{%
\subsubsection{A mixed amendment needs two revenue calculations}\label{a-mixed-amendment-needs-two-revenue-calculations}}

The modification record retains the original allocation/performance, enforceable change, remaining distinctness findings, amended allocations, faithful measures of progress and new receipts. AASB~15.21(c) \citep[paras.~18--21,39--45,73--90]{aasb15dated} requires consistent application of 21(a)/(b) to the affected promises; it does not supply a universal price split. The following exact case deliberately removes discount-allocation ambiguity.

A contract includes construction100 and four distinct inspection visits5 each, at their evidenced standalone prices. Inspections concern another operating line and can occur before the new construction finishes. The original expected construction cost80; actual eligible cost40 and faithful cost-to-cost progress give50\% and revenue50. One visit adds revenue5 and cash expense2. The120 price was prepaid. The approved amendment adds integrated construction40 and two visits10; supported standalone allocations therefore become140 and30. New total construction cost100 makes cumulative progress40\%, so target construction revenue56 yields catch-up6. The first visit's5 remains; five remaining distinct visits share25 prospectively. The new50 receipt itself earns no revenue.

With opening cash/equity100, two additional visits by year-end yield revenue71, expense46, profit25, cash224 and contract liability99. The indirect operating bridge25+99=124 matches170 received−46 paid. Complete construction for cash60 and three visits for6 next year: revenue84+15=99, profit33 and cash/equity158. Lifetime profit58=price170−construction100−visit costs12. Keeping old progress would recognise70 for construction; applying40\% to the entire170 would wrongly recognise distinct future visits. Missing allocation/progress or enforceability blocks the numeric transition. The one-year-or-less financing expedient under63 is explicitly elected for this example, not assumed for long prepaid contracts.

\hypertarget{unfinished-production-carries-component-specific-completion}{%
\subsubsection{Unfinished production carries component-specific completion}\label{unfinished-production-carries-component-specific-completion}}

The production record adds started/completed/scrapped quantities, material insertion point, conversion stages, opening WIP by component, observed progress, actual eligible costs, normal capacity, fixed/variable overhead classification, NRV, and the costing-policy version. Equivalent units are a derived policy justified by this homogeneous process; IAS~2 \citep[IAS~2.12--18]{eu20231803} prescribes rational cost allocation and the normal-capacity constraint, not this exact spreadsheet convention.

Start100 units with no opening WIP or abnormal waste. Finish80;20 are half converted, with all material already inserted. Material1000 is spread over100 material units; variable conversion450 over90 equivalent conversion units. Fixed cash overhead600 has admitted normal capacity120 equivalent units, so its rate is5, assigned450 and expensed idle150. Material rate10 plus variable5 plus fixed5 gives completed unit cost20. Completed cost1600 and WIP300 exhaust assigned cost1900: WIP has material200 and conversion100. Sell50 finished units at30: revenue1500, cost1000 and idle expense150 give profit350; finished600 plus WIP300 remain. Opening cash/equity3000 leaves cash2450 and equity3350.

In the next production month, the remaining10 equivalent conversion units require variable50 and fixed50; independently supplied normal capacity10 admits both. Add100 to WIP300 and transfer completed400. The remaining50 finished units then sell1500 at total cost1000. Cumulative cash/equity3850 and profit850 follow; no WIP remains. A different opening WIP, spoilage policy or completion pattern needs its own allocation. Reusing100 as conversion units overstates completion; capitalising all600 fixed cost conceals150 idle expense. NRV is separately assessed for both unfinished and completed stock, including costs to complete.

\hypertarget{joint-products-and-delayed-supplier-payment}{%
\subsubsection{Joint products and delayed supplier payment}\label{joint-products-and-delayed-supplier-payment}}

JointThe production record retains inseparable joint costs, split-off quantities, separate subsequent costs, allocation-basis evidence, selling values/dates and NRV. For joint cost300, A60 at split-off price4 and B40 at9 give values240/360. The selected rational and consistent relative-sales method assigns120/180, or unit costs2/4.5. Selling30A and10B gives revenue210, cost105 and profit105. Opening cash500 leaves cash410, unsold A60/B135 and equity605. Weighting by mass instead changes the admitted policy; equal weights are not a missing-value fix.

Separately, production300 with an independently established immaterial byproduct of NRV20 allocates main inventory280 and byproduct20. Sale of that byproduct for20 removes its20 carrying amount; it earns no margin here. Opening cash500 ends cash220+main280=equity500. Retaining main300 as well as byproduct20 invents cost; deducting20 without recording the unsold byproduct loses an asset.

The supplier record names the original creditor, normal-credit purchase price, actual terms, financing assessment, initial liability, effective rate, payment schedule and actual payments. Normal cash price100 and a fixed payment121 in two years imply10\% effective interest under the supplied terms. Record inventory and payable100; interest10 then11 raises the liability121. Ordinary inventory is not a qualifying asset in this case. Sell for cash150 in year1: opening cash/equity100 becomes cash250, payable110 and equity140. Year2 interest11 and payment121 leave cash/equity129; lifetime profit29. Under the explicit historical IAS7 election, both original supplier payment and interest are operating cash. Year1 indirect40+payable increase110=150; year2−11−110=−121. A substituted finance provider or altered arrangement requires reclassification analysis.

\hypertarget{uncertain-tax-has-its-own-measurement-decision}{%
\subsubsection{Uncertain tax has its own measurement decision}\label{uncertain-tax-has-its-own-measurement-decision}}

The prior tax dossier already specifies tax bases, gross DTA/DTL, recoverability, rate changes and lease interaction; the error/estimate bridge retains equity origins. Uncertain tax treatment was genuinely deferred. Under the historical EU IFRIC23.6--14 method, The uncertain-tax record adds treatment grouping, filing position, authority/examination scope, full-information acceptance finding, candidate outcomes, predictive-measure choice, consistent current/deferred bases, reassessment event and disclosure obligations. Assume examination with full knowledge; do not multiply exposure by an estimated chance of avoiding examination.

Synthetic law ultimately deducts equipment100, with the timing uncertain. Book depreciation20 leaves PPE80; cash operating revenue100 makes pretax profit80. Opening cash/equity500, purchase100 and cash revenue100 leave cash500. The enacted rate25\% is supplied, all tax amounts remain payable, and all origins are profit. If full deduction100 is probably accepted, current tax0 and DTL20 follow. If outcomes100/20 are binary with full acceptance probability1/5, and most-likely amount better predicts resolution, use deduction20: current20 and DTL0. In a third branch, deductions100/40/20 at1/5,3/10,1/2 and a supported expected-value choice give42. Taxable profit58 and tax base58 give current14.5 and DTL5.5. All branches have total tax20, profit60 and equity560; they differ in the timing liability. These equal totals arise from this particular full-deduction law.

If new evidence immediately resolves the expected-value branch at deduction40, increase current tax14.5→15 and reduce DTL5.5→5. The two0.5 expense/benefit entries net to zero. Pay15: cash485+PPE80=DTL5+equity560. A later resolution must also include intervening depreciation and other facts; silence alone is not new evidence of acceptance. Real law, probability and predictive suitability remain required. The printed hypothetical rate or distribution cannot fill them.

\hypertarget{combining-the-new-operating-contributions}{%
\subsubsection{Combining the new operating contributions}\label{combining-the-new-operating-contributions}}

Combine only the five explicitly compatible EU year1 contributions: WIP, joint products, byproduct, supplier financing and lifetime restoration. Give each independent lots/contracts/assets but one reporting entity, currency and period; sum their openings, without intercompany elimination or transferring AU cases. Opening cash4200, PPE120, provision20 and equity4300 become cash3430, inventory1375, PPE71, liabilities143 and pretax equity4733. Pretax profit433 reconciles equity4300→4733; total assets4876 equal liabilities143+equity4733. Operating cash movement−770 is350+105+0+40−62 profits, plus depreciation60, less inventory increase1375, plus supplier increase110, plus noncash unwind2. The11 restoration estimate adds both PPE and provision and is excluded from cash. This is a pretax operating slice; new tax effects are unresolved, so a complete after-tax report cannot be issued from it. The separate uncertain-tax example is not silently grafted onto this different law/fact pattern.
